\begin{register}{H}{LPPERI\_CLK\_EN\_REG}{0x0000}\label{regdesc:LPPERICLKENREG}
\regfield{(reserved)}{1}{31}{0}%
\regfield{LPPERI\_EFUSE\_CK\_EN}{1}{30}{{1}}%
\regfield{(reserved)}{1}{29}{1}%
\regfield{LPPERI\_LP\_EXT\_I2C\_CK\_EN}{1}{28}{{1}}%
\regfield{(reserved)}{3}{25}{111}%
\regfield{LPPERI\_RNG\_CK\_EN}{1}{24}{{1}}%
\regfield{(reserved)}{24}{0}{000000000000000000000000}%
\reglabel{Reset}\regnewline%
\begin{regdesc}\begin{reglist}
\label{fielddesc:LPPERIRNGCKEN}\item [LPPERI\_RNG\_CK\_EN] Configures whether to enable the RNG clock.\\
0: Disable\\
1: Enable\\
When \nameref{mod:ecdsa} is active, the RNG clock is forced on regardless of this field.\\
(R/W)
\label{fielddesc:LPPERILPEXTI2CCKEN}\item [LPPERI\_LP\_EXT\_I2C\_CK\_EN] Configures whether to enable the LP external I2C clock.\\
0: Disable\\
1: Enable\\
(R/W)
\label{fielddesc:LPPERIEFUSECKEN}\item [LPPERI\_EFUSE\_CK\_EN] Configures whether to enable the eFuse controller clock.\\
0: Disable\\
1: Enable\\
(R/W)
\end{reglist}\end{regdesc}
\end{register}


\begin{register}{H}{LPPERI\_RESET\_EN\_REG}{0x0004}\label{regdesc:LPPERIRESETENREG}
\regfield{(reserved)}{1}{31}{0}%
\regfield{LPPERI\_EFUSE\_RESET\_EN}{1}{30}{{0}}%
\regfield{(reserved)}{5}{25}{00000}%
\regfield{LPPERI\_LP\_RNG\_RESET\_EN}{1}{24}{{0}}%
\regfield{(reserved)}{24}{0}{000000000000000000000000}%
\reglabel{Reset}\regnewline%
\begin{regdesc}\begin{reglist}
\label{fielddesc:LPPERILPRNGRESETEN}\item [LPPERI\_LP\_RNG\_RESET\_EN] Configures whether to assert reset to the LP RNG.\\
0: Release reset\\
1: Assert reset\\
(R/W)
\label{fielddesc:LPPERIEFUSERESETEN}\item [LPPERI\_EFUSE\_RESET\_EN] Configures whether to assert reset to the eFuse controller.\\
0: Release reset\\
1: Assert reset\\
(R/W)
\end{reglist}\end{regdesc}
\end{register}


\begin{register}{H}{LPPERI\_RNG\_DATA\_REG}{0x0008}\label{regdesc:LPPERIRNGDATAREG}
\regfield{LPPERI\_RND\_DATA}{32}{0}{{0}}%
\reglabel{Reset}\regnewline%
\begin{regdesc}\begin{reglist}
\label{fielddesc:LPPERIRNDDATA}\item [LPPERI\_RND\_DATA] Represents asynchronous RNG data before synchronization. For cryptographic-quality random data, read \hyperref[regdesc:LPPERIRNGDATASYNCREG]{LPPERI\_RNG\_DATA\_SYNC\_REG} instead. (RO)
\end{reglist}\end{regdesc}
\end{register}


\begin{register}{H}{LPPERI\_CPU\_REG}{0x000C}\label{regdesc:LPPERICPUREG}
\regfield{LPPERI\_LPCORE\_DBGM\_UNAVALIABLE}{1}{31}{{1}}%
\regfield{(reserved)}{31}{0}{0000000000000000000000000000000}%
\reglabel{Reset}\regnewline%
\begin{regdesc}\begin{reglist}
\label{fielddesc:LPPERILPCOREDBGMUNAVALIABLE}\item [LPPERI\_LPCORE\_DBGM\_UNAVALIABLE] Configures the LP CPU Debug Mode state.\\
0: Available\\
1: Unavailable\\
(R/W)
\end{reglist}\end{regdesc}
\end{register}


\begin{register}{H}{LPPERI\_BUS\_TIMEOUT\_REG}{0x0010}\label{regdesc:LPPERIBUSTIMEOUTREG}
\regfield{LPPERI\_LP\_PERI\_TIMEOUT\_PROTECT\_EN}{1}{31}{{1}}%
\regfield{LPPERI\_LP\_PERI\_TIMEOUT\_INT\_CLEAR}{1}{30}{{0}}%
\regfield{LPPERI\_LP\_PERI\_TIMEOUT\_THRES}{16}{14}{{0xffff}}%
\regfield{(reserved)}{14}{0}{00000000000000}%
\reglabel{Reset}\regnewline%
\begin{regdesc}\begin{reglist}
\label{fielddesc:LPPERILPPERITIMEOUTTHRES}\item [LPPERI\_LP\_PERI\_TIMEOUT\_THRES] Configures the timeout threshold for accesses to LP peripheral registers.\\
Measurement unit: number of clock cycles in the timeout clock domain.\\
(R/W)
\label{fielddesc:LPPERILPPERITIMEOUTINTCLEAR}\item [LPPERI\_LP\_PERI\_TIMEOUT\_INT\_CLEAR] Write 1 to clear the LP peripheral bus timeout interrupt. (WT)
\label{fielddesc:LPPERILPPERITIMEOUTPROTECTEN}\item [LPPERI\_LP\_PERI\_TIMEOUT\_PROTECT\_EN] Configures whether to enable timeout protection for accesses to LP peripheral registers.\\
0: Disable\\
1: Enable\\
(R/W)
\end{reglist}\end{regdesc}
\end{register}


\begin{register}{H}{LPPERI\_BUS\_TIMEOUT\_ADDR\_REG}{0x0014}\label{regdesc:LPPERIBUSTIMEOUTADDRREG}
\regfield{LPPERI\_LP\_PERI\_TIMEOUT\_ADDR}{32}{0}{{0}}%
\reglabel{Reset}\regnewline%
\begin{regdesc}\begin{reglist}
\label{fielddesc:LPPERILPPERITIMEOUTADDR}\item [LPPERI\_LP\_PERI\_TIMEOUT\_ADDR] Represents the address of the access that triggered an LP peripheral bus timeout. (RO)
\end{reglist}\end{regdesc}
\end{register}


\begin{register}{H}{LPPERI\_BUS\_TIMEOUT\_UID\_REG}{0x0018}\label{regdesc:LPPERIBUSTIMEOUTUIDREG}
\regfield{(reserved)}{25}{7}{0000000000000000000000000}%
\regfield{LPPERI\_LP\_PERI\_TIMEOUT\_UID}{7}{0}{{0x0}}%
\reglabel{Reset}\regnewline%
\begin{regdesc}\begin{reglist}
\label{fielddesc:LPPERILPPERITIMEOUTUID}\item [LPPERI\_LP\_PERI\_TIMEOUT\_UID] Represents the master ID and permission captured when an LP peripheral bus timeout occurs. Bits [4:0] are the master ID and bits [6:5] are the master permission. See Chapter \ref{mod:permctrl} \textit{\nameref{mod:permctrl}}. Cleared when the timeout interrupt is cleared. (RO)
\end{reglist}\end{regdesc}
\end{register}


\begin{register}{H}{LPPERI\_MEM\_CTRL\_REG}{0x001C}\label{regdesc:LPPERIMEMCTRLREG}
\regfield{LPPERI\_UART\_MEM\_FORCE\_PU}{1}{31}{{1}}%
\regfield{LPPERI\_UART\_MEM\_FORCE\_PD}{1}{30}{{0}}%
\regfield{LPPERI\_UART\_WAKEUP\_EN}{1}{29}{{0}}%
\regfield{(reserved)}{27}{2}{000000000000000000000000000}%
\regfield{LPPERI\_UART\_WAKEUP\_FLAG}{1}{1}{{0}}%
\regfield{LPPERI\_UART\_WAKEUP\_FLAG\_CLR}{1}{0}{{0}}%
\reglabel{Reset}\regnewline%
\begin{regdesc}\begin{reglist}
\label{fielddesc:LPPERIUARTWAKEUPFLAGCLR}\item [LPPERI\_UART\_WAKEUP\_FLAG\_CLR] Write 1 to clear LPPERI\_UART\_WAKEUP\_FLAG. (WT)
\label{fielddesc:LPPERIUARTWAKEUPFLAG}\item [LPPERI\_UART\_WAKEUP\_FLAG] Represents whether an LP UART wakeup has occurred.\\
0: No LP UART wakeup occurred\\
1: LP UART wakeup occurred\\
(R/WTC/SS)
\label{fielddesc:LPPERIUARTWAKEUPEN}\item [LPPERI\_UART\_WAKEUP\_EN] Configures whether to enable the LP UART wakeup function.\\
0: Enable\\
1: Disable\\
(R/W)
\label{fielddesc:LPPERIUARTMEMFORCEPD}\item [LPPERI\_UART\_MEM\_FORCE\_PD] Configures whether to force the LP UART memory to power down.\\
0: Do not force power down\\
1: Force power down\\
(R/W)
\label{fielddesc:LPPERIUARTMEMFORCEPU}\item [LPPERI\_UART\_MEM\_FORCE\_PU] Configures whether to force the LP UART memory to power up.\\
0: Do not force power up\\
1: Force power up\\
(R/W)
\end{reglist}\end{regdesc}
\end{register}


\begin{register}{H}{LPPERI\_INTERRUPT\_SOURCE\_REG}{0x0020}\label{regdesc:LPPERIINTERRUPTSOURCEREG}
\regfield{(reserved)}{26}{6}{00000000000000000000000000}%
\regfield{LPPERI\_LP\_INTERRUPT\_SOURCE}{6}{0}{{0x0}}%
\reglabel{Reset}\regnewline%
\begin{regdesc}\begin{reglist}
\label{fielddesc:LPPERILPINTERRUPTSOURCE}\item [LPPERI\_LP\_INTERRUPT\_SOURCE] Represents the LP interrupt source.\\
Bit 5: PMU\_LP\_INT\\
Bit 4: Reserved\\
Bit 3: RTC\_TIMER\_LP\_INT\\
Bit 2: LP\_UART\_INT\\
Bit 1: LP\_I2C\_INT\\
Bit 0: LP\_IO\_INT\\
(RO)
\end{reglist}\end{regdesc}
\end{register}


\begin{register}{H}{LPPERI\_RNG\_CFG\_REG}{0x0024}\label{regdesc:LPPERIRNGCFGREG}
\regfield{(reserved)}{20}{12}{00000000000000000000}%
\regfield{LPPERI\_RTC\_TIMER\_EN}{2}{10}{{3}}%
\regfield{(reserved)}{9}{1}{111111111}%
\regfield{LPPERI\_RNG\_SAMPLE\_ENABLE}{1}{0}{{0}}%
\reglabel{Reset}\regnewline%
\begin{regdesc}\begin{reglist}
\label{fielddesc:LPPERIRNGSAMPLEENABLE}\item [LPPERI\_RNG\_SAMPLE\_ENABLE] Configures whether to enable BUF\_CHAIN (ring-oscillator seed path).\\
0: Disable\\
1: Enable\\
(R/W)
\label{fielddesc:LPPERIRTCTIMEREN}\item [LPPERI\_RTC\_TIMER\_EN] Bit [0]: Configures whether the RTC timer is used as a noise source input to CRC.\\
0: Disable\\
1: Enable\\
Bit [1]: Configures whether to superimpose the RTC timer output on the random number output by CRC to increase randomness.\\
0: Disable\\
1: Enable\\
(R/W)
\end{reglist}\end{regdesc}
\end{register}


\begin{register}{H}{LPPERI\_RNG\_DATA\_SYNC\_REG}{0x0028}\label{regdesc:LPPERIRNGDATASYNCREG}
\regfield{LPPERI\_RND\_SYNC\_DATA}{32}{0}{{0}}%
\reglabel{Reset}\regnewline%
\begin{regdesc}\begin{reglist}
\label{fielddesc:LPPERIRNDSYNCDATA}\item [LPPERI\_RND\_SYNC\_DATA] Represents synchronized true random data from the RNG. (RO)
\end{reglist}\end{regdesc}
\end{register}


\begin{register}{H}{LPPERI\_DATE\_REG}{0x03FC}\label{regdesc:LPPERIDATEREG}
\regfield{LPPERI\_CLK\_EN}{1}{31}{{0}}%
\regfield{LPPERI\_LPPERI\_DATE}{31}{0}{{0x2312170}}%
\reglabel{Reset}\regnewline%
\begin{regdesc}\begin{reglist}
\label{fielddesc:LPPERILPPERIDATE}\item [LPPERI\_LPPERI\_DATE] Version control register. (R/W)
\label{fielddesc:LPPERICLKEN}\item [LPPERI\_CLK\_EN] Configures whether to enable the diagnostic clock for accesses to this register block.\\
0: Disable\\
1: Enable\\
Use only for diagnostics. Keep this field cleared to 0 in normal operation.\\
(R/W)
\end{reglist}\end{regdesc}
\end{register}
